% !TeX root = ../main.tex

\chapter{Introduction}
\label{ch:introduction}

\section{Motivation}
\label{sec:intro:motivation}

Edge computing has shifted general-purpose computation from centralized data centers toward devices and points-of-presence located close to data sources and end users. Two pressures shape the workloads that run in this regime. The first is a tight per-request response-time budget; tens of milliseconds are routinely required for interactive content delivery, real-time inference, and request-handling at the network edge. The second is a tight per-host resource budget. Edge nodes are typically provisioned with substantially less CPU, memory, and energy headroom than data-center servers. The workload mix is dominated by short-lived, stateless functions that frequently start and exit.

WebAssembly (Wasm)~\cite{haas2017wasm,wasm-spec} has emerged as the deployment vehicle for these workloads. Its compact binary format, language-agnostic toolchain, sandboxed execution model, and predictable resource semantics make it well suited to multi-tenant environments in which untrusted code must be loaded, executed, and torn down at high frequency. Production systems including Fastly Compute~\cite{fastly-compute}, Cloudflare Workers~\cite{cloudflare-workers}, and Fermyon Spin~\cite{spin-docs} embed Wasm runtimes for exactly this reason. The class of workloads these systems target is one for which cold-start latency is as critical as steady-state throughput~\cite{gackstatter2022serverless}: a serverless function invoked once and discarded pays the full cost of compilation before producing any user-visible output. This thesis uses \emph{cold-start latency} to mean the time from request arrival at the runtime to the start of execution of the first machine-code instruction of the requested wasm function. The definition assumes no ready-to-run compiled instance of the module is already resident on the host, and excludes the wasm function's own execution time, which is reported separately as steady-state work.

These cloud-native runtimes share a common deployment model. Each incoming request is dispatched to a freshly-instantiated Wasm sandbox that loads, executes, and tears down within the request handler's lifetime. The model enables strong tenant isolation without the kernel-level overhead of containers or virtual machines, and the language-agnostic toolchain accepts modules compiled from C, Rust, Go, and other source languages without runtime-specific build adaptations. The compile-time profile of the underlying runtime therefore bounds the response-time floor that the platform can expose to the user. Production deployments hide a portion of this cost through ahead-of-time compilation caches and pre-warmed sandbox pools, but neither workaround eliminates the per-tenant cold path that emerges whenever a previously-unseen module is dispatched on a node. The structural limitation is that single-tier optimizing compilation forces every invocation through the same compile pipeline, regardless of whether the request will execute the module for one microsecond or for one minute.

WasmEdge~\cite{wasmedge-docs} is an open-source Wasm runtime developed under the Cloud Native Computing Foundation and adopted in production by ByteDance, Microsoft, Sony, and others. Its compile path uses LLVM end-to-end for both ahead-of-time (AOT) and just-in-time (JIT) compilation. LLVM produces highly optimized machine code. Its optimization pipeline includes mid-level inlining, loop transformations, vectorization, and global register allocation. The result is throughput close to native C compilation on computationally intensive workloads. These optimizations come at a cost. LLVM's compile-time profile scales poorly with module size, and its memory footprint during compilation is substantial. On the resource-constrained hosts and short-lived functions that characterize edge deployments, paying this cost on every invocation forces a structural trade-off: the runtime cannot deliver a low cold-start latency without sacrificing peak code quality.

A second observation complicates the single-tier design. Within a typical Wasm module, the distribution of execution time across functions is highly skewed: a small fraction of functions accounts for the bulk of run-time, while the remainder execute infrequently or not at all. Compiling every function with LLVM-grade optimization spends most of the compile budget on code that will not run hot. Cold or rarely-executed code receives optimization it does not need, and the compile time imposed on hot code delays the program from reaching its hot path.

This pattern is well-established in the design of high-performance language runtimes. Modern JavaScript engines---V8~\cite{v8-wasm}, JavaScriptCore~\cite{webkit-wasm}, SpiderMonkey~\cite{firefox-tiering}---all employ tiered compilation. A fast baseline tier produces native code for every function with minimal optimization. One or more optimizing tiers then selectively recompile hot code with progressively heavier optimization, guided by run-time profile information. The baseline tier eliminates the cold-start gap; the optimizing tier delivers steady-state performance comparable to a single-tier optimizing compiler. Wasm runtimes built atop these engines, including V8's Liftoff and JavaScriptCore's BBQ tier, follow the same architecture~\cite{v8-wasm,webkit-wasm}. WasmEdge does not. The runtime currently lacks a baseline JIT, lacks a profile-driven promotion mechanism, and lacks a mid-execution code-replacement pathway. The cold-start bottleneck observed in production is a direct consequence of this architectural gap.

\section{Problem Statement}
\label{sec:intro:problem}

This thesis addresses the absence of a tiered compilation infrastructure in WasmEdge. The goal is to introduce an architecture that decouples cold-start latency from peak code quality, while preserving the steady-state performance that LLVM-compiled code achieves on hot workloads. Three concrete sub-problems constitute this goal.

\paragraph{Baseline tier.} The runtime requires a compiler that produces native code for every wasm function at low cost at module load time. The compiler must integrate with WasmEdge's existing runtime contracts---linear memory layout, table dispatch, host-call boundaries, trap handling, and the function-instance store. It must lower the full set of MVP WebAssembly instructions plus the reference-type, bulk-memory, and table extensions used by sightglass and real workloads. It must produce code of sufficient quality that the resulting steady-state performance is acceptable until tier-2 compilation completes. A lightweight SSA-based code generator with linear-scan register allocation is an established design point that meets these constraints.

\paragraph{Optimizing tier.} The runtime requires a path that recompiles hot wasm functions with LLVM-grade optimization, but only after profile information has identified them as worth the compile cost. Reimplementing WasmEdge's existing wasm-to-LLVM lowering is undesirable: the LLVM frontend is mature, well-tested, and already encodes the module-level lowering logic that selective recompilation also requires. Two challenges arise. The first is to drive the existing LLVM frontend with input that contains only the hot functions, while preserving the call-graph structure on which the frontend's call lowering depends. The second is to bridge the resulting LLVM-compiled code into the baseline tier's calling convention so that mixed-tier dispatch is correct. The promotion mechanism must run asynchronously to user code so that LLVM's compile latency is hidden from the request-handling path, and the published code must replace the baseline-tier entry atomically with respect to concurrent callers.

\paragraph{In-flight migration.} Function-entry promotion alone is insufficient for an important class of workloads. Consider a function called once that spends the bulk of its execution inside a single long-running loop. This is the shape of cryptographic kernels, parsing loops, and many serverless request handlers. The function never re-enters its prologue after the counter saturates, so the entry-table swap delivers no speedup on the in-flight invocation. Closing this gap requires a mechanism that detects when an executing tier-1 frame has crossed a hotness threshold inside a loop, migrates the frame's wasm-level state into a tier-2 continuation entry, and transfers control mid-execution. The mechanism must coexist with function-entry promotion, must impose negligible overhead in the steady state once it has fired, and must preserve wasm semantics across the transition.

Each sub-problem presents system-level challenges that go beyond the compilation algorithm itself. The list includes the choice of profile representation and storage, concurrent publication of native code, lifetime management of code caches that outlive their compiling thread, and the interaction between tier-1 and tier-2 calling conventions on the indirect-call paths that Wasm modules use heavily. The remainder of this thesis presents a design that resolves each of these constraints, an implementation in WasmEdge, and an evaluation of the resulting system.

\section{Contributions}
\label{sec:intro:contributions}

This thesis makes the following contributions.

\begin{enumerate}
  \item \textbf{A baseline JIT tier (Tier-1) for WasmEdge.} The proposed tier integrates a lightweight Sea-of-Nodes JIT backend as a compile-every-function pathway invoked at module instantiation. The backend runs its standard optimization pipeline. Profile counters embedded in the compiled bodies drive subsequent tier-2 promotion.

  \item \textbf{A selective LLVM-frontend recompilation tier (Tier-2).} The proposed tier reuses the existing WasmEdge LLVM frontend by synthesizing per-batch mini-modules that contain only hot functions. A tier-up heuristic selects which functions enter each batch. Three ABI thunks bridge the two tiers' calling conventions.

  \item \textbf{An on-stack replacement (OSR) mechanism for tier-1 to tier-2 migration.} The proposed mechanism instruments outermost-loop back-edges in tier-1 code to detect hot loops. The mechanism then migrates the executing tier-1 frame into a tier-2 continuation entry mid-loop. The migration closes the gap on workloads where function-entry promotion alone delivers no speedup.

  \item \textbf{Empirical evaluation.} The thesis evaluates the architecture on the sightglass benchmark suite, with whole-module LLVM JIT as the steady-state-throughput baseline. The chapter thoroughly analyzes the results per workload shape and isolates the contribution of each tier-2 mechanism.
\end{enumerate}

The architecture is implemented in WasmEdge. The tier-up mechanism is concise. The evaluation shows it is also effective on the edge and serverless workloads it targets. Cold-start latency drops by up to one order of magnitude on the real-world modules typical of these deployments. Steady-state throughput stays close to whole-module LLVM JIT. To the best of the author's knowledge, this combination is not documented in the browser-engine or server-side WebAssembly literature in this exact form.

\section{Thesis Organization}
\label{sec:intro:organization}

The remainder of this thesis is organized as follows. Chapter~\ref{ch:background} reviews the WebAssembly execution model, the WasmEdge runtime, the structural cost of LLVM compilation, the dstogov/ir library, the principles of tiered compilation, and prior work on tiered compilation in JavaScript and Wasm runtimes, on-stack replacement, and lightweight JIT libraries. Chapter~\ref{ch:design} presents the design of the proposed system: the architecture overview, the runtime context and calling conventions, the tier-1 baseline JIT, the tier-2 selective recompilation pipeline, and the on-stack replacement mechanism. Chapter~\ref{ch:evaluation} reports the evaluation: methodology, results across the three components, cold-start latency, and sensitivity sweeps. Chapter~\ref{ch:conclusion} concludes the thesis and discusses limitations and future work.
